\section{An elementary determinant estimate for quadratic characters}
\label{sz-section}

This section isolates the part of the suggested auxiliary literature that can
be proved with particularly little machinery. The argument originates in the
OpenAI proof of the exclusion of Landau--Siegel zeros and its subsequent
simplifications, including Kerr's earlier interpolation argument
\cite{sz-kerr}. The rectangular form below follows the use of derivatives
in Green's argument, as explained by Tao, and the reduction modulo the largest
inert prime used by Pescador and van Dijk
\cite{sz-openai,sz-tao,sz-pescador,sz-vandijk}.
We give all of the determinant argument, with an explicit finite inequality
and a slightly sharper bound for the determinant. No result about automorphic
representations is used in this section.

For a nonsquare integer $d$, set $A=\mathbb Z[\sqrt d]$ and write
$\overline{u+v\sqrt d}=u-v\sqrt d$. An odd prime $p$ will be called inert
here if the Legendre symbol $(d/p)$ is $-1$. This definition makes sense even
when $A$ is not the maximal order. It implies
\begin{equation}\label{sz-frobenius}
 A/pA\simeq\mathbb F_p[T]/(T^2-d)\simeq\mathbb F_{p^2},
 \qquad \alpha^p\equiv\bar\alpha\pmod {pA}\quad(\alpha\in A).
\end{equation}
Indeed, $T^2-d$ is irreducible modulo $p$, and Euler's criterion gives
$(\sqrt d)^p=-\sqrt d$ in that field.

\begin{lemma}[Interpolation in characteristic $p$]\label{sz-interpolation}
Let $k$ be a field of characteristic $p$, let $1\le s\le p$, and let
$a_1,\ldots,a_p$ be distinct elements of $k$. The map
\[
 F\longmapsto\bigl((\partial_z^jF)(a_i,a_i^p)\bigr)_{
          1\le i\le p,\ 0\le j<s}
\]
from the polynomials with $\deg_z F<p$ and $\deg_w F<s$ to $k^{ps}$
is an isomorphism.
\end{lemma}
\begin{proof}
The source and target have dimension $ps$, so injectivity suffices. Put
$G(z)=F(z,z^p)$. The monomials $z^rw^t$ become the distinct monomials
$z^{r+pt}$, for $0\le r<p$ and $0\le t<s$. Thus $F\ne0$ implies
$G\ne0$ and $\deg G<ps$. Since $(z^p)'=0$, the chain rule gives
$G^{(j)}(a_i)=(\partial_z^jF)(a_i,a_i^p)$ for every $j$. Vanishing of the
displayed evaluations forces $G$ to have a zero of order at least $s$ at
each $a_i$: the coefficients of $(z-a_i)^j$, for $j<s$, are
$G^{(j)}(a_i)/j!$, and these factorials are invertible because $j<p$.
The distinct factors $(z-a_i)^s$ then give $\deg G\ge ps$, a contradiction.
\end{proof}

\begin{proposition}[A finite determinant inequality]\label{sz-finite}
Let $R$ be an odd inert prime, let $1\le S\le R$, and let
$\alpha_1,\ldots,\alpha_R\in A$ have distinct reductions modulo $RA$.
Suppose that both complex embeddings of every $\alpha_i$ have modulus
at most $H$, where $H\ge1$. Put $N=RS$ and, for each odd inert prime $p$, put
\begin{align}
 E_p&=\sum_{r=0}^{R-1}\sum_{s=0}^{S-1}
        \min\{\lfloor r/p\rfloor,S-1-s\}\label{sz-E}\\
    &=\sum_{j=1}^{S-1}(S-j)\max\{R-jp,0\}.\nonumber
\end{align}
Then
\begin{equation}\label{sz-finite-ineq}
 2\sum_{\substack{p\ \mathrm{odd}\\ (d/p)=-1}} E_p\log p
 \le N\log N+N(R-1)\log H+N(S-1)\log R.
\end{equation}
The sum is finite, since $E_p=0$ for $p\ge R$.
\end{proposition}
\begin{proof}
Form the $N\times N$ matrix
\[
 B_{(r,s),(i,j)}=(\partial_z^j z^rw^s)(\alpha_i,\bar\alpha_i),
 \quad 0\le r<R,\quad 0\le s<S,\quad1\le i\le R,\quad0\le j<S.
\]
Its determinant $\Delta$ belongs to $A$. Reducing modulo $RA$,
\eqref{sz-frobenius} and Lemma~\ref{sz-interpolation} show that $B$ is
invertible. In particular $\Delta\ne0$. Consequently
$M=|\Delta\bar\Delta|$ is a positive rational integer.

Fix an odd inert prime $p$ and set
$h_{r,s}=\min\{\lfloor r/p\rfloor,S-1-s\}$. Replace the row
corresponding to $z^rw^s$ by the evaluation row of
\[
 z^{r-ph_{r,s}}(z^p-w)^{h_{r,s}}w^s.
\]
The leading term is $z^rw^s$. Every other term has smaller $z$-degree
and still has $z$-degree below $R$ and $w$-degree below $S$.
Ordering the rows by their $z$-degree shows that the change of basis is
triangular over $\mathbb Z$ with diagonal entries $1$, so it preserves
$\Delta$.

Every ordinary $z$-derivative of $z^p-w$, evaluated at
$(\alpha_i,\bar\alpha_i)$, lies in $pA$. For derivative order $0$ this
is \eqref{sz-frobenius}; for orders $1\le j\le p$ the coefficient
$p(p-1)\cdots(p-j+1)$ is divisible by $p$; for higher orders the derivative
is zero. In the product rule each of the $h_{r,s}$ factors therefore
contributes a factor $p$, even if some factors are differentiated.
The transformed row lies in $p^{h_{r,s}}A^N$. Factoring these powers
from the rows gives $\Delta\in p^{E_p}A$, and taking its norm gives
$p^{2E_p}\mid M$. Distinct primes may be multiplied, whence
\begin{equation}\label{sz-norm-lower}
 \log M\ge2\sum_p E_p\log p.
\end{equation}
The second formula for $E_p$ counts, for each $j\ge1$, the pairs
$(r,s)$ with $r\ge jp$ and $s\le S-1-j$.

We bound the determinant separately under each complex embedding.
An entry with $j\le r$ has modulus at most
$R^jH^{r+s-j}$, and an entry with $j>r$ is zero. Divide row $(r,s)$
by $H^{r+s}$ and column $(i,j)$ by $(R/H)^j$. Every resulting entry
has modulus at most $1$, so Hadamard's inequality bounds its determinant
by $N^{N/2}$. Moreover,
\[
 \sum_{r,s}(r+s)-\sum_{i,j}j=\frac{N(R-1)}2,
 \qquad \sum_{i,j}j=\frac{N(S-1)}2.
\]
Thus
\[
 |\Delta|\le N^{N/2}H^{N(R-1)/2}R^{N(S-1)/2}
\]
under either embedding. Multiplying the bounds for the two embeddings,
taking logarithms, and using \eqref{sz-norm-lower} proves the result.
\end{proof}

\begin{remark}[A number-field version of the finite inequality]
Proposition~\ref{sz-finite} only needs the following data: a number field
$K$ of degree $e$, an order $\mathcal A\subset K$, an automorphism
$\tau$ of $\mathcal A$, and elements $\alpha_1,\ldots,\alpha_R$.
For every prime $p$ to be used in the lower bound, assume
$\alpha_i^p-\tau(\alpha_i)\in p\mathcal A$ for every $i$.
Require the same Frobenius congruence at the auxiliary prime $R$,
and assume also that the images of the
$\alpha_i$ in some residue field of characteristic $R$ are distinct.
If all embeddings of the $\alpha_i$ and $\tau(\alpha_i)$ have modulus
at most $H$, the identical inequality \eqref{sz-finite-ineq} holds, with
the sum restricted to these primes. To prove this, replace
$\bar\alpha_i$ by $\tau(\alpha_i)$ in the matrix. Reduction to the
chosen residue field again makes the determinant nonzero. The row changes
give $\Delta\in p^{E_p}\mathcal A$, so
$p^{eE_p}$ divides the nonzero integer $N_{K/\mathbb Q}(\Delta)$.
The archimedean calculation gives the $e$th power of the bound for one
embedding. Multiplying the logarithmic comparison by $2/e$ gives exactly
\eqref{sz-finite-ineq}. No unramifiedness assertion is implicit in this
version; the stated congruences and residue-field hypothesis are required.
\end{remark}

\begin{remark}[Ordinary derivatives are essential]
One cannot replace every $\partial_z^j$ by the divided derivative
$\partial_z^j/j!$ in the divisibility proof. At order $p$ the divided
derivative of $z^p-w$ is $1$, which is not divisible by $p$.
The restriction $S\le R$ concerns the auxiliary prime proving nonvanishing;
the other inert primes can be much smaller than $S$. Thus invertibility
of factorials at $R$ does not justify dividing by them at all those primes.
\end{remark}

Here are the elementary prime estimates needed to pass from the finite
inequality to a uniform statement. All implied constants below are absolute
and effective.

\begin{lemma}\label{sz-mertens}
For $x\ge2$ and $y\ge2$,
\[
 \vartheta(x):=\sum_{p\le x}\log p\ll x,
 \qquad \sum_{p\le x}\frac{\log p}{p}=\log x+O(1),
 \qquad \sum_{p\ge y}\frac{\log p}{p^2}\ll y^{-1}.
\]
\end{lemma}
\begin{proof}
Every prime between $n$ and $2n$ divides $\binom{2n}{n}$, so
$\vartheta(2n)-\vartheta(n)\le2n\log2$. Summation on dyadic intervals
gives $\vartheta(x)\ll x$. Hence, writing $\Lambda$ for the von Mangoldt
function,
\[
 \psi(x):=\sum_{m\le x}\Lambda(m)
 =\sum_{k\ge1}\vartheta(x^{1/k})
 \ll x+\sqrt x\log(2x)\ll x.
\]
For an integer $n\ge2$, prime factorization of $n!$ gives
\[
 \log(n!)=\sum_{m\le n}\Lambda(m)\lfloor n/m\rfloor
 =n\sum_{m\le n}\frac{\Lambda(m)}m+O(\psi(n)).
\]
Integral comparison gives $\log(n!)=n\log n+O(n)$, so
$\sum_{m\le n}\Lambda(m)/m=\log n+O(1)$. The terms with $m=p^k$,
$k\ge2$, have bounded total, because
$\sum_{p,k\ge2}(\log p)/p^k\le\sum_{n\ge2}(\log n)/(n(n-1))<\infty$.
This proves the second estimate, also for real $x$ by taking its integer
part. The last estimate follows by partial summation from
$\vartheta(t)\ll t$.
\end{proof}

\begin{theorem}[Uniform bound for inert primes]\label{sz-inert}
For every nonsquare integer $d$ and every $X\ge2$,
\begin{equation}\label{sz-inert-bound}
 I_d(X):=\sum_{\substack{p\le X\ p\ {
m odd}\ (d/p)=-1}}
            \frac{\log p}{p}
 \le\frac12\log X+\frac14\log|d|+\log\log(3X)+O(1).
\end{equation}
The implied constant is independent of both $d$ and $X$.
\end{theorem}
\begin{proof}
If there is no odd inert prime below $X$, the assertion is immediate.
Otherwise let $R$ be the largest such prime. Then $I_d(X)=I_d(R)$.
If $R<e^3$, the result is immediate with an absolute constant. If
$R\le\sqrt{|d|}$, Lemma~\ref{sz-mertens} gives
\[
 I_d(R)\le\log R+O(1)
 \le\tfrac12\log R+\tfrac14\log|d|+O(1).
\]
We may therefore assume $R\ge e^3$ and $R>\sqrt{|d|}$.

Set $L=\log R$, $S=\lceil R/L\rceil$, and $N=RS$. Then $1\le S<R$.
Consider the quadratic integers $u+v\sqrt d$ with
\[
 0\le u<\sqrt R\,|d|^{1/4},
 \qquad0\le v<\sqrt R\,|d|^{-1/4},\qquad u,v\in\mathbb Z.
\]
There are at least $R$ such integers because the two side lengths have
product $R$. Each coordinate lies in $[0,R)$, so all reductions are distinct
in $A/RA$. Select any $R$ of them. Both embeddings have modulus below
$H=2\sqrt R\,|d|^{1/4}<2R$. Apply Proposition~\ref{sz-finite},
and let $M$ be the positive determinant norm constructed in its proof.

For $a\in\mathbb Z_{\ge0}$ one has
$\sum_{s<S}\min\{a,S-1-s\}\ge Sa-a^2$.
For $a<S$ this follows by summing the arithmetic progression; for $a\ge S$
the right side is nonpositive. Consequently
\begin{align*}
 E_p&\ge S\sum_{r<R}\lfloor r/p\rfloor-
              \sum_{r<R}\lfloor r/p\rfloor^2\\
 &\ge\frac{SR(R-1)}{2p}-SR-\frac{R^3}{3p^2}.
\end{align*}
Sum this inequality only over inert primes with $L\le p\le R$.
By Lemma~\ref{sz-mertens}, and because $S\ge R/L$,
\[
 2\sum_pE_p\log p
 \ge NR\sum_{\substack{L\le p\le R\\(d/p)=-1}}
           \frac{\log p}{p}-O(NR).
\]
Indeed, replacing $R-1$ by $R$ costs $O(N\log R)$, the floor errors
cost $O(N\vartheta(R))$, and the quadratic error is
$O(R^3/L)=O(NR)$.
On the other hand, \eqref{sz-finite-ineq} gives
\[
 \frac{\log M}{NR}
 \le\frac{\log N}{R}+\left(1-\frac1R\right)\log H
       +\frac{S-1}{R}\log R
 \le\log H+O(1).
\]
Here $(S-1)\log R/R\le1$ and $N\le R^2$.
Restore the primes below $L$, whose contribution is at most
$\log L+O(1)$ by Lemma~\ref{sz-mertens}. The result is
\[
 I_d(R)\le\log H+\log L+O(1)
 =\tfrac12\log R+\tfrac14\log|d|+\log\log R+O(1).
\]
Finally $R\le X$ proves \eqref{sz-inert-bound}.
\end{proof}

The parameter choice can be checked directly instead of guessed. The same
calculation with any integer $1\le S\le R/2$, and cutoff $R/S$ for the
inert primes, gives
\[
 I_d(R)\le\log H+\frac SR\log R+\log(R/S)+O(1).
\]
Writing $u=S/R$, the varying part of this upper bound is
$u\log R-\log u$. Its derivative is $\log R-1/u$, so its minimum occurs
at $u=1/\log R$. This recovers the choice above. Increasing the number of
derivatives arbitrarily does not improve this estimate: it increases the
term $u\log R$. This calculation describes a limitation of this particular
determinant bound, and is not a claim of optimality among all possible
arithmetic constructions.

\subsection{The analytic consequence and its exact scope}

A nonprincipal real primitive Dirichlet character is quadratic:
there is a fundamental discriminant $D$ with $|D|=q$ for which
$\chi(n)=(D/n)$, the Kronecker symbol. A fundamental discriminant is
the discriminant of a quadratic number field. In particular, there is
a squarefree $d\ne1$, $|d|\le q$, such that $\chi(p)=(d/p)$ at every
odd prime. Thus Theorem~\ref{sz-inert} applies to the inert primes
of each real primitive character.

The following use of complex analysis is explicit. We use the classical
analytic continuation, functional equation, and Hadamard product for a
primitive Dirichlet $L$-function. If $a=(1-\chi(-1))/2$, these give, for
real $\sigma>1$ and real primitive nonprincipal $\chi$,
\begin{equation}\label{sz-hadamard}
 \sum_\rho\Re\frac1{\sigma-\rho}
 =\frac12\log(q/\pi)
   +\frac12\frac{\Gamma'}{\Gamma}\left(\frac{\sigma+a}{2}\right)
   +\frac{L'}{L}(\sigma,\chi).
\end{equation}
The zeros here are the nontrivial zeros, with multiplicity; all summands
are nonnegative and their sum converges. The same identity holds at
$\sigma=1$ when $L(1,\chi)\ne0$. These are classical analytic inputs,
not consequences of the determinant calculation. They are also the inputs
used in the last step of \cite{sz-openai,sz-pescador,sz-vandijk}.

\begin{corollary}\label{sz-zero-gap}
Every real zero $\beta\in(0,1)$ of a real primitive nonprincipal character
of conductor $q$ satisfies
\[
 (1-\beta)\log q\ge1-O\left(\frac{\log\log(3q)}{\log q}\right).
\]
In particular there is an effective absolute constant $c>0$ such that
$(1-\beta)\log q\ge c$ for all these characters and zeros.
\end{corollary}
\begin{proof}
For $0<\eta\le1$, partial summation in Theorem~\ref{sz-inert}, including
the harmless prime $2$, gives
\begin{equation}\label{sz-laplace}
 \sum_{\chi(p)=-1}\frac{\log p}{p^{1+\eta}}
 \le\frac1{2\eta}+\frac14\log q+O(\log(2/\eta)).
\end{equation}
For completeness, the sum is
$\eta\int_1^\infty S_-(t)t^{-1-\eta}\,dt$, where
$S_-(t)=\sum_{p\le t,\chi(p)=-1}(\log p)/p$.
The integrals of $1$ and $\log t$ against this probability measure are
$1$ and $1/\eta$. With $u=\log t$, concavity of the logarithm bounds
the integral of $\log(\log3+u)$ by $\log(\log3+1/\eta)$.

The Euler products, with all prime powers of exponent at least two
contributing $O(1)$, give
\begin{align*}
 \frac{L'}{L}(1+\eta,\chi)
 &=2\sum_{\chi(p)=-1}\frac{\log p}{p^{1+\eta}}
     -\sum_p\frac{\log p}{p^{1+\eta}}
     +\sum_{p\mid q}\frac{\log p}{p^{1+\eta}}+O(1)\\
 &\le\frac12\log q+O(\log(2/\eta)+\log\log(3q)).
\end{align*}
Here $\sum_p(\log p)/p^{1+\eta}=1/\eta+O(1)$ follows by partial
summation from Lemma~\ref{sz-mertens}. Also
$\sum_{p\mid q}(\log p)/p\ll\log\log(3q)$: split at $p=\log q$,
use Lemma~\ref{sz-mertens} below that point, and use
$\sum_{p\mid q}\log p\le\log q$ above it.

Retain only $\rho=\beta$ in \eqref{sz-hadamard}. The gamma term is
bounded for $1\le\sigma\le2$, uniformly in the parity, and hence
\[
 \frac1{1-\beta+\eta}
 \le\log q+O(\log(2/\eta)+\log\log(3q)).
\]
For large $q$, take $\eta=(\log q)^{-2}$ and invert the positive
right side. This gives the asserted asymptotic inequality and a positive
uniform constant for all sufficiently large $q$.

To see that the extension to the finitely many smaller conductors is
effective, use the classical theorem $L(1,\chi)\ne0$. Periodicity gives
$|\sum_{u<n\le v}\chi(n)|\le q/2$ and hence an effective tail bound
$q/(2N)$ for $\sum_{n>N}\chi(n)/n$. Computing partial sums therefore
eventually certifies a positive lower bound for $|L(1,\chi)|$.
Summation by parts also gives
$|L'(s,\chi)|\le 2q/e$ for $1/2\le s\le1$: the total variation of
$(\log t)t^{-s}$ on $[1,\infty)$ is $2/(es)$, and the partial sums
are bounded by $q/2$. The mean value inequality supplies a positive
zero-free interval at $1$. Taking the minimum over the finite list and
also over $\tfrac12\log q$ for zeros below $1/2$ gives an effective $c$.
\end{proof}

The conclusion is a logarithmic gap from $1$, not a half-plane with boundary
$1/2$. Neither a nonreal zero with real part $3/4$ nor a real zero bounded
away from $1$ contradicts the inequalities just proved. For example, the
bound obtained from \eqref{sz-hadamard} at a hypothetical real zero
$\beta=3/4$ compares a number tending to $4$ with a quantity of size
$\log q$. It cannot exclude such a zero as $q$ grows. This observation
identifies the exact limitation of this method; it does not assert that
any Dirichlet $L$-function has such a zero.

